
%%%%%%%%%%%%%%%%%%%%%%%%%%%%%%%%%%%%%%%%
%           main body
%%%%%%%%%%%%%%%%%%%%%%%%%%%%%%%%%%%%%%%%


%-----------------------------------
%	TITLE PAGE
%-----------------------------------

\title[Linear Algebra]{\LARGE \bfseries 第二部分\quad 线性代数} % The short title appears at the bottom of every slide, the full title is only on the title page
\author[Exercises] % Your institution as it will appear on the bottom of every slide, may be shorthand to save space
{\Large \bfseries 习题}
% \author{Singular Value Decomposition} % Your name
% \date{\today} % Date, can be changed to a custom date
\date{}

%------------------------------------------------

\begin{document}

%------------------------------------------------

\begin{frame}

\titlepage % Print the title page as the first slide

\end{frame}

%------------------------------------------------

% \begin{frame}
% \frametitle{内容提要} % Table of contents slide, comment this block out to remove it
% \tableofcontents % Throughout your presentation, if you choose to use \section{} and \subsection{} commands, these will automatically be printed on this slide as an overview of your presentation
% \end{frame}

%-----------------------------------
%	PRESENTATION SLIDES
%-----------------------------------

%------------------------------------------------

\begin{frame}

{\color{red} \scriptsize 带$^*$号的题目可以不做}

\vspace{1em}

1. 用高斯消元法解下列线性方程组
\begin{itemize}
\item[(1)] $\displaystyle \systeme{-x_2 -x_3 + x_4 = 0, x_1 + x_2 + 3x_3 - x_4 = 0, -2x_1 + 3x_2 - x_4 = 0, -x_1 - 4x_3 - 4x_4 = 0}$
\item[(2)] $\displaystyle \systeme{x_1 + x_2 + x_3 = 0, x_2 + x_3 = -2, -x_1 - 2x_3 = -2}$
\end{itemize}

\end{frame}

%------------------------------------------------

\begin{frame}

2. 计算下列方阵的行列式
% \begin{itemize}
% \item[(1)] $\displaystyle $
% \item[(2)] $\displaystyle $
% \end{itemize}

(1) $\begin{pmatrix} 1 & 1 \\ 1 & 1 \end{pmatrix}$ \quad (2) $\begin{pmatrix} \cos\theta & -\sin\theta \\ \sin\theta & \cos\theta \end{pmatrix}$ \quad (3) $\begin{pmatrix} \cos\theta & 0 & -\sin\theta \\ 0 & 1 & 0 \\ \sin\theta & 0 & \cos\theta \end{pmatrix}$ \newline
\vspace{0.5em}
(4) $\begin{pmatrix} 1 & 2 & -2 & 1 \\ -1 & -2 & 4 & -2 \\ 0 & 1 & 2 & -1 \\ 4 & 3 & -3 & 5 \end{pmatrix}$ \quad (5) $\begin{pmatrix} 1 & 1 & 1 & 1 \\ 1 & -1 & 2 & 3 \\ 1 & 1 & 4 & 9 \\ 1 & -1 & 8 & 27 \end{pmatrix}$

\vspace{1em}

3$^*$. 设$\left\{ \begin{array}{rcl} L_1: \alpha x + \beta y + \gamma = 0 \\ L_2: \gamma x + \alpha y + \beta = 0 \\ L_3: \beta x + \gamma y + \alpha = 0\end{array}\right.$是$\mathbb{R}^2$平面上三条不同的直线。如果$L_1, L_2, L_3$交于一点，试证$\alpha + \beta + \gamma = 0$

\end{frame}

%------------------------------------------------

\begin{frame}

4. 计算矩阵的乘积
\begin{itemize}
\item[(1)] 设$A = \begin{pmatrix} a_1 \\ a_2 \end{pmatrix}, B = \begin{pmatrix} b_1 & b_2 \end{pmatrix}$, 分别计算$AB$与$BA$
\vspace{0.5em}
\item[(2)] 设$A = \begin{pmatrix} 1 & 1 \\ -1 & -1 \end{pmatrix}, B = \begin{pmatrix} 1 & -1 \\ -1 & 1 \end{pmatrix}$, 分别计算$AB$与$BA$
\vspace{0.5em}
\item[(3)] 设$A = \begin{pmatrix} 1 & 2 \\ 2 & 4 \end{pmatrix}, B = \begin{pmatrix} -1 & 3 \\ -2 & 1 \end{pmatrix}, C = \begin{pmatrix} -7 & 1 \\ 1 & 2 \end{pmatrix}$，分别计算$AB$与$AC$
\vspace{0.5em}
\item[(4)] 设$A = \begin{pmatrix} a_1 \\ a_2 \\ a_3 \end{pmatrix}, B = \begin{pmatrix} b_1 & b_2 & b_3 \end{pmatrix}$，求$(AB)^{100}$
\end{itemize}

\end{frame}

%------------------------------------------------

\begin{frame}

5. 求下列方阵的逆
\newline \vspace{0.5em}
(1) $\begin{pmatrix} \cos\varphi & -\sin\varphi \\ \sin\varphi & \cos\varphi \end{pmatrix}$ \quad (2) $\begin{pmatrix} a & b \\ c & d \end{pmatrix}$, 其中$ad-bc\neq 0$
\newline \vspace{0.5em}
(3) $\begin{pmatrix} 2 & 0 & 7 \\ -1 & 4 & 5 \\ 3 & 1 & 2 \end{pmatrix}$ \qquad\quad (4) $\begin{pmatrix} 1 & 1 & 1 & 1 \\ 1 & 1 & -1 & -1 \\ 1 & -1 & 1 & -1 \\ 1 & -1 & -1 & 1 \end{pmatrix}$

\vspace{1em}

6. 设$n$阶方阵$A$满足$A^k = 0$, 求证$I_n-A$可逆，且它的逆为
\[(I_n-A)^{-1} = I_n + A + A^2 + \cdots + A^{k-1}\]

\end{frame}

%------------------------------------------------

\begin{frame}

7$^*$. 设$A$为$n$阶可逆阵。
\begin{itemize}
\item[(1)] 设$\mathbf{u}, \mathbf{v}$是$n$维列向量，且$1 + \mathbf{v}^T A^{-1} \mathbf{u} \neq 0$，求证
\[(A+\mathbf{u}\mathbf{v}^T)^{-1} = A^{-1} - \frac{A^{-1} \mathbf{u} \mathbf{v}^T A^{-1}}{1 + \mathbf{v}^T A^{-1} \mathbf{u}}\]
\item[(2)] 设$U$是$n\times k$矩阵，$V$是$k\times n$矩阵，$C$是$k\times k$矩阵，且$ C^{-1}+VA^{-1}U$可逆，求证
\[ \left(A+UCV \right)^{-1} = A^{-1} - A^{-1}U \left(C^{-1}+VA^{-1}U \right)^{-1} VA^{-1}.\]
\end{itemize}

\end{frame}

%------------------------------------------------

\begin{frame}

8. 判断下列子集是否为给定线性空间$\mathbb{R}^3$的子空间\\
\begin{itemize}
\item[(1)] $V_1 = \{(x_1,2,0) \in \mathbb{R}^3\}$
\vspace{0.4em}
\item[(2)] $V_2 = \{(x_1,0,x_3) \in \mathbb{R}^3\}$
\vspace{0.4em}
\item[(3)] $V_3 = \{(x_1,x_2,x_3) \in \mathbb{R}^3 \ |\ x_1-3x_2+2x_3=0\}$
\vspace{0.4em}
\item[(4)] $V_4 = \{(x_1,x_2,x_3) \in \mathbb{R}^3 \ |\ x_1-3x_2+2x_3=1\}$
\vspace{0.4em}
\item[(5)] $V_5 = \{(x_1,x_2,x_3) \in \mathbb{R}^3 \ |\ \frac{x_1}{2}=\frac{x_2-4}{3}=\frac{x_3-3}{4}\}$
\vspace{0.4em}
\item[(6)] $V_6 = \{(x_1,x_2,x_3) \in \mathbb{R}^3 \ |\ x_1=x_3\mbox{且}x_1+x_2+x_3=0\}$
\end{itemize}

\end{frame}

%------------------------------------------------

\begin{frame}

9. 设$\alpha = (1,-2,2,0)^T, \beta = (3,2,1,-2)^T$为$\mathbb{R}^4$中的两个向量
\begin{itemize}
\item[(1)] 分别求$\alpha, \beta$的长度，二者的内积以及夹角
\vspace{0.4em}
\item[(2)] 设$V = \operatorname{span}\{\alpha, \beta\}$, 求$V$在$\mathbb{R}^4$中的正交补$V^{\perp}$
\end{itemize}

\vspace{1em}

10. 用Gram-Schmidt正交化程序把$\mathbb{R}^4$的基
\[(1,1,1,1)^T,(1,0,1,1)^T,(1,1,0,1)^T,(1,1,1,0)^T\]
化为一组标准正交基

\vspace{1em}

11. 将实对称阵$A = \begin{pmatrix} 0 & 1 & 1 & -1 \\ 1 & 0 & -1 & 1 \\ 1 & -1 & 0 & 1 \\ - 1 & 1 & 1 & 0 \end{pmatrix}$对角化

\end{frame}

%------------------------------------------------

\begin{frame}

12. 考虑矩阵$A = \begin{pmatrix} 1 & 1 & 0 \\ 0 & 1 & 1 \\ 0 & 0 & 5 \end{pmatrix}$与矩阵$B = \begin{pmatrix} 1 & 0 & 0 \\ 0 & 1 & 0 \\ 0 & 0 & 5 \end{pmatrix}$. 问
\begin{itemize}
\item[(1)] $A,B$是否相抵?
\item[(2)] $A,B$是否相似?
\end{itemize}

\vspace{1em}

13. 假设在某省，每年有比例$p$的农村居民移居城镇，有比例为$q$的城镇居民移居到农村。假设该省的人口总数不变，
且上述人口迁移的规律不变。将$n$ 年后农村人口和城镇人口占总人口的比例分别记为$x_n$ 和$y_n$ （$x_n+y_n=1$）。
\begin{itemize}
\item[(1)] 求关系式$\begin{pmatrix} x_{n+1} \\ y_{n+1} \end{pmatrix} = A \begin{pmatrix} x_n \\ y_n \end{pmatrix}$ 中的矩阵$A$
\item[(2)] 设目前农村与城镇人口相等，即$\begin{pmatrix} x_0 \\ y_0 \end{pmatrix} = \begin{pmatrix} 0.5 \\ 0.5 \end{pmatrix}$,
求$\begin{pmatrix} x_{n} \\ y_{n} \end{pmatrix}$
\end{itemize}

\end{frame}

%------------------------------------------------

\begin{frame}

14. 判别下列二次型是否正定
\begin{itemize}
\item[(1)] $10x_1^2 + 8x_1x_2 + 24x_1x_3 + 2x_2^2 - 28x_2x_3 + x_3^2$
\vspace{0.4em}
\item[(2)] $\sum\limits_{i=1}^n x_i^2 + \sum\limits_{1\le i<j\le n} x_ix_j$
\end{itemize}

\vspace{1em}

15. 考虑二次型
\[Q(x_1,x_2,x_3) = \lambda(x_1^2+x_2^2+x_3^2) + 2x_1x_2 + 2x_1x_3 - 2x_2x_3\]
问$\lambda$分别取什么值的时候，$Q$是正定，以及负定的

\end{frame}

%------------------------------------------------

\begin{frame}

16. 求以下矩阵的Moore-Penrose伪逆$A^+$

\vspace{0.5em}

(1) $(a,b,c)$ \quad (2) $\begin{pmatrix} a \\ b \\ c \end{pmatrix}$ \quad (3) $\begin{pmatrix} a & 0 & 0 \\ 0 & b & 0 \\ 0 & 0 & 0 \end{pmatrix}$ \quad (4) $\begin{pmatrix} a & b \\ 2a & 2b \end{pmatrix}$

\vspace{0.5em}

(5) $\begin{pmatrix} 3 & 0 \\ 4 & 5 \end{pmatrix}$ \quad (6) $\begin{pmatrix} 0 & 1 & & \\ & \ddots & \ddots & \\ & & \ddots & 1 \\ & & & 0 \end{pmatrix}$ \quad (7) $\begin{pmatrix} 0 & \cdots & 0 & 1 \\ \vdots & \ddots & & 0 \\ \vdots & & \ddots & \vdots \\ 0 & \cdots & \cdots & 0 \end{pmatrix}$

\end{frame}

%------------------------------------------------

\begin{frame}

17. 求下列线性方程组的最小二乘解
\begin{itemize}
\item[(1)] $\begin{pmatrix} 1 & -2 & 4 \\ 1 & -1 & 1 \\ 1 & 0 & 0 \\ 1 & 1 & 1 \\ 1 & 2 & 4 \end{pmatrix} \begin{pmatrix} x_1 \\ x_2 \\ x_3 \end{pmatrix} = \begin{pmatrix} 0 \\ 1 \\ 2 \\ 1 \\ 0 \end{pmatrix}$
\item[(2)] $\begin{pmatrix} 1 & 1 & 1 \\ 1 & 1 & 1 \end{pmatrix} \begin{pmatrix} x_1 \\ x_2 \\ x_3 \end{pmatrix} = \begin{pmatrix} 1 \\ 2 \end{pmatrix}$
\end{itemize}

\vspace{1em}

18$^*$. 叙述并证明复矩阵情况的奇异值分解定理

\end{frame}

%------------------------------------------------

% \begin{frame}

% \begin{block}{Vandermonde行列式的计算}
% \begin{enumerate}
% % \addtocounter{enumi}{0}
% \item 把第$n-1$行乘以$-a_1$加到第$n$行，再把第$n-2$行乘以$-a_1$加到第$n-1$行。按照这种顺序依次向上，直到用第$2$行乘以$-a_1$加到第$1$行。由行列式性质，$V_n$的值不变。
% \[
% V_n = \det \begin{pmatrix}
% 1 & 1 & \cdots & 1 \\
% 0 & a_2-a_1 & \cdots & a_n-a_1 \\
% \vdots & \vdots & & \vdots \\
% 0 & a_2^{n-2}(a_2-a_1) & \cdots & a_n^{n-2}(a_n-a_1) \\
% \end{pmatrix}
% \]
% \end{enumerate}
% \end{block}

% \end{frame}

%------------------------------------------------
% % closing page
% \begin{frame}
% \HUGE{\centerline{\bfseries {\color{gray} The} {\color{zkblue} End}}}

% \pause

% \vspace{1.2em}

% \Large{\centerline{\bfseries {\color{gray}Thank} \quad {\color{zkblue} You}}}

% \end{frame}

%----------------------------------------------------------------------------------------

\end{document}
